%% Elevark fra mattelab.no – laget automatisk av bygg.py. Rediger gjerne til eget bruk.
%% Kompiler med lualatex.
\documentclass[11pt,a4paper]{article}
\usepackage[norsk]{babel}
\usepackage[a4paper,margin=2.2cm,top=2cm,bottom=2.4cm]{geometry}
\usepackage{fontspec}
\usepackage{unicode-math}
\setmainfont{TeX Gyre Pagella}
\setsansfont{TeX Gyre Heros}[Scale=MatchLowercase]
\setmathfont{TeX Gyre Pagella Math}
\usepackage{xcolor}
\definecolor{blekk}{HTML}{16213A}
\definecolor{dempet}{HTML}{566077}
\definecolor{rute}{HTML}{CFC9B8}
\definecolor{m1}{HTML}{2F4AB0}
\definecolor{m2}{HTML}{99266D}
\usepackage[most]{tcolorbox}
\usepackage{qrcode}
\usepackage{needspace}
\usepackage{fancyhdr}
\usepackage[hidelinks]{hyperref}
\pdfvariable suppressoptionalinfo \numexpr 1+2+4+8+16+32+64+128+256+512\relax
\setlength{\parindent}{0pt}
\setlength{\parskip}{0.4em}
\pagestyle{fancy}
\fancyhf{}
\renewcommand{\headrulewidth}{0pt}
\fancyfoot[L]{\sffamily\footnotesize\color{dempet}mattelab.no · Sekant som blir tangent · elevark uten fasit}
\fancyfoot[R]{\sffamily\footnotesize\color{dempet}Side \thepage}

\newcommand{\fase}[2]{\tcbox[on line, colframe=#2, colback=white, coltext=#2, boxrule=0.6pt, arc=2pt,
  boxsep=0pt, left=3pt, right=3pt, top=1.5pt, bottom=1.5pt, fontupper=\sffamily\bfseries\footnotesize]{#1}}
\newcommand{\steg}[4]{%
  \par\noindent{\color{rute}\rule{\linewidth}{0.4pt}}\par\nopagebreak
  \noindent\begin{minipage}[t]{2.1cm}\vspace{0pt}\fase{#1}{#2}\end{minipage}%
  \begin{minipage}[t]{\dimexpr\linewidth-2.1cm\relax}\vspace{0pt}#3\end{minipage}\par
  \vspace{#4}}
\newcommand{\del}[4]{%
  \par\noindent{\color{rute}\rule{\linewidth}{0.4pt}}\par\nopagebreak
  \noindent\begin{minipage}[t]{1.1cm}\vspace{0pt}\fase{#1}{#2}\end{minipage}%
  \begin{minipage}[t]{\dimexpr\linewidth-1.1cm\relax}\vspace{0pt}#3\end{minipage}\par
  \vspace{#4}}
\newcommand{\oppgave}[2]{\par\Needspace{12\baselineskip}\vspace{0.6em}%
  {\fontsize{14}{17}\selectfont\bfseries{\color{m1}#1}\enspace #2}\par\vspace{0.2em}}

\begin{document}
\begin{minipage}[t]{0.72\linewidth}\vspace{0pt}
{\sffamily\small\color{dempet}Elevark · 1T og R1 · mattelab.no}\par\vspace{0.3em}
{\fontsize{24}{28}\selectfont\bfseries Sekant som blir tangent}\par\vspace{1em}
{\sffamily\small Navn: \rule{4.6cm}{0.4pt}\quad Klasse: \rule{1.5cm}{0.4pt}\quad Dato: \rule{1.9cm}{0.4pt}}
\end{minipage}\hfill
\begin{minipage}[t]{0.22\linewidth}\vspace{0pt}\raggedleft
\qrcode[height=2.3cm]{https://mattelab.no/utforskere/sekant-tangent.html}\par\vspace{0.2em}
{\sffamily\scriptsize\color{dempet}Utforskeren}
\end{minipage}

\vspace{1em}
{\sffamily\small\color{dempet}Utforskeren ligger på \url{https://mattelab.no/utforskere/sekant-tangent.html}}\par
Hver oppgave har tre steg. \textbf{Gjett} først, uten å bruke utforskeren. \textbf{Test} så i utforskeren. \textbf{Forklar} til slutt med egne ord hvorfor det blir slik.

\oppgave{1}{Gjennomsnittlig vekstfart}
\steg{Gjett}{m1}{Hva er den gjennomsnittlige vekstfarten til $f(x) = x^2$ fra $x = 1$ til $x = 3$?}{2.2cm}
\steg{Test}{m2}{Velg $f(x) = x^2$ med $x_0 = 1$ og $h = 2$. Les av $\Delta x$, $\Delta y$ og stigningstallet til sekanten.}{2.6cm}
\steg{Forklar}{blekk}{Hvorfor er den gjennomsnittlige vekstfarten det samme som stigningstallet til sekanten?}{3.4cm}

\oppgave{2}{La $h$ gå mot null}
\steg{Gjett}{m1}{Fortsatt med $f(x) = x^2$ og $x_0 = 1$: Hva tror du skjer med stigningstallet til sekanten når $h$ blir $1$, $0{,}1$, $0{,}01$ og $0{,}001$?}{2.2cm}
\steg{Test}{m2}{Sett $x_0 = 1$ og $h = 1$. Trykk «La h gå mot 0», og se på tabellen. Prøv også negative $h$.}{2.6cm}
\steg{Forklar}{blekk}{Vis at stigningstallet til sekanten er $2 + h$. Hva skjer når $h \to 0$? Hvorfor kan vi ikke bare sette $h = 0$?}{3.4cm}

\oppgave{3}{Den deriverte som funksjon}
\steg{Gjett}{m1}{For $f(x) = x^2$: Hva tror du stigningstallet til tangenten er i $x = -1$, $x = 0$ og $x = 2$?}{2.2cm}
\steg{Test}{m2}{Sett $x_0 = -1$ og $h = 0{,}001$. Dra $x_0$ til $0$ og så til $2$. Følg med på tangentens stigningstall.}{2.6cm}
\steg{Forklar}{blekk}{Vis at $\frac{(x + h)^2 - x^2}{h} = 2x + h$, og bruk det til å finne $f'(x)$.}{3.4cm}

\oppgave{4}{Vannrett tangent}
\steg{Gjett}{m1}{Grafen til $f(x) = x^3 - 3x$ har et toppunkt og et bunnpunkt. I hvilke $x$-verdier tror du tangenten er vannrett?}{2.2cm}
\steg{Test}{m2}{Velg $f(x) = x^3 - 3x$. Dra $x_0$ til du finner punktene der tangentens stigningstall er $0$.}{2.6cm}
\steg{Forklar}{blekk}{Den deriverte er $f'(x) = 3x^2 - 3$. Løs $f'(x) = 0$. Hvorfor er tangenten vannrett i et toppunkt eller bunnpunkt?}{3.4cm}

\oppgave{5}{Når finnes ikke den deriverte?}
\steg{Gjett}{m1}{Grafen til $f(x) = |x|$ har en knekk i $x = 0$. Hva tror du skjer med sekanten når $h$ går mot $0$ fra høyre? Fra venstre?}{2.2cm}
\steg{Test}{m2}{Velg $f(x) = |x|$ med $x_0 = 0$. Se på tabellen for positive og negative $h$. Velg så $f(x) = \sqrt{x}$ med $x_0 = 0$, og la $h$ gå mot $0$.}{2.6cm}
\steg{Forklar}{blekk}{Forklar hvorfor $f'(0)$ ikke finnes for noen av de to funksjonene. Er grunnen den samme?}{3.4cm}

\end{document}
